\documentclass[aps,prd,10pt,superscriptaddress,nofootinbib,showkeys,amsmath,amssymb]{revtex4-2}

\usepackage{graphicx}
\usepackage{bm}
\usepackage{booktabs}
\usepackage{array}
\usepackage{amsthm}
\usepackage{enumerate}
\usepackage{hyperref}
\usepackage{lineno}

\hypersetup{
  colorlinks=true,
  linkcolor=blue,
  citecolor=blue,
  urlcolor=blue
}

\newtheorem{hypothesis}{Hypothesis}
\newtheorem{prediction}{Prediction}
\newtheorem{criterion}{Criterion}

\begin{document}

\title{Paper XI: Dark Matter, Dark Energy, Black Holes, and the Blancken-Layer Synthesis: A Constraint-Driven Flux Interpretation of Cosmological Anomalies}

\author{Steve Bico Mujjabi, MD}
\email{msbico@gmail.com}
\homepage{https://github.com/bampita-bico/CDFD-Part-I-Release}
\thanks{Independent researcher; ORCID: \href{https://orcid.org/0009-0001-0556-5516}{0009-0001-0556-5516}.}
\affiliation{VuraLabs, Kampala, Uganda}

\date{July 2026}

\begin{abstract}
Paper XI develops the speculative cosmological bridge of Constraint-Driven Flux Dynamics (CDFD). It asks whether dark matter phenomenology, dark energy, black-hole horizons, vacuum self-consistency, and the proposed sub-Planckian Blancken layer can be treated as related manifestations of constrained transport rather than unrelated mysteries. The manuscript is deliberately status-aware: the lepton and vacuum equations in Papers I--X contain formal derivations internal to the model, while the dark-sector and Blancken-layer claims remain hypotheses until they generate observational discriminants beyond qualitative analogy.

\medskip
\noindent In the extended notation used across the series, the vacuum limit is written
as $\Psi_s = (\Phi/C)\cdot S\cdot M_s$. This paper treats that expression as a
hypothesis generator for cosmological discriminants, not as established
cosmological evidence. The public CDFD Runtime implements the same
state grammar for reusable simulations; the paper-local supplementary script
remains the audited reproduction layer for the figures and tables in this Part I
archive.
\end{abstract}

\keywords{Constraint-Driven Field Theory, Constraint-Driven Flux Dynamics, adaptive operating ratio, vacuum structure, Mujjabi tests, reproducible physics, dark matter, dark energy, black holes}

\maketitle
\linenumbers

\section*{Project Context}
This manuscript constitutes Part I (Paper 11) of the multi-part series
\textit{Constraint-Driven Flux Dynamics (CDFD)}. The underlying mathematical
architecture builds directly upon the concepts established in Paper 10 of this
series (Steve Bico Mujjabi, \textit{``CDFD Part I: Fundamental Physics ---
Paper 10: The Vacuum Equation of State: Deriving Torus Knot Self-Consistency
from the Public CDFD Balance Law $\Psi_s = \Phi / C$''}, prepared for APS submission and archived in the Part I REVTeX release).

\section*{Symbol Conventions}

\noindent\textbf{CDFD physical-vacuum symbol convention.}
\begin{center}
{\small
\begin{tabular}{p{0.16\linewidth}p{0.33\linewidth}p{0.41\linewidth}}
\hline
Symbol & Part I meaning & Physical reading \\
\hline
$J$ or $\Phi$ & Driving flux or throughput & Transport current, field throughput, circulation drive, or generic flux; a subscript fixes the exact channel \\
$C$ & Active constraint or capacity burden & Vacuum transport capacity, geometric confinement, stiffness, boundary resistance, or load-bearing limit \\
$S$ & Surface responsiveness & Local ability of the vacuum or modeled surface to reconfigure under drive \\
$M_s$ & Structural memory & Retained geometry, phase history, stiffness history, or accumulated constraint state \\
$\Psi_s$ & Adaptive operating ratio $(J/C)\cdot S\cdot M_s$ or $(\Phi/C)\cdot S\cdot M_s$ & Local bookkeeping gauge for constrained operation; it is not a new universal constant \\
$R,a,\chi$ & Ring radius, core radius, and aspect ratio $R/a$ & Geometry used to connect vortex structure to dimensionless electromagnetic scales \\
$\rho_0,c,\kappa$ & Vacuum density scale, propagation speed, and transport stiffness & Medium parameters used in stability, pressure, and self-consistency calculations \\
\hline
\end{tabular}
}
\end{center}

\noindent\textbf{Greek-symbol convention.}
\begin{center}
{\small
\begin{tabular}{p{0.16\linewidth}p{0.33\linewidth}p{0.41\linewidth}}
\hline
Symbol & Local role & Interpretation rule \\
\hline
$\alpha$ & Fine-structure constant unless a local equation states otherwise & In the Part I physics papers, un-subscripted $\alpha$ means the electromagnetic fine-structure constant \\
$\beta,\gamma$ & Local response, stiffness, coupling, or correction coefficients & Equation-local coefficients; subscripts bind them to the named mode, channel, or approximation \\
$\theta,\phi,\Phi_c$ & Phase, orientation, or chirality angles & Used for torus-knot orientation, vacuum phase, or chiral phase geometry as defined in the nearby equation \\
$\delta,\epsilon$ & Perturbation, residual, tolerance, or small-offset scales & Local stability margins, numerical residuals, or experimental tolerances \\
$\lambda,\mu,\nu$ & Rate, mass, mode, or frequency labels & Meaning is fixed by the equation, subscript, and physics context where the symbol appears \\
$\rho,\sigma,\tau$ & Density, spread/noise, or time-scale terms & Local density, variance, stress, relaxation, or transport-memory quantities \\
$\Omega,\Lambda$ & State/domain space or integrated measure where defined & Reserved for explicitly defined state spaces, spectra, or synthesis-level quantities \\
\hline
\end{tabular}
}
\end{center}

\section{Motivation}

Modern cosmology contains several successful but conceptually separate descriptions.
Dark matter is inferred from galaxy dynamics, cluster lensing, and precision
cosmology~\cite{RubinFord1970,Clowe2006,Planck2018}. Dark energy is inferred
from supernova acceleration and the cosmological parameter ladder
~\cite{Riess1998,Perlmutter1999,Planck2018}. Black holes are described by
general relativity, horizon mechanics, and semiclassical thermodynamics
~\cite{BardeenCarterHawking1973,Hawking1975}. Particle masses are treated
within quantum field theory. CDFD does not replace these mature descriptions.
Its proposed contribution is a unifying transport ontology:
\begin{equation}
    \Psi_{s,\mathrm{vac}}=\frac{\Phi_{\mathrm{vac}}}{C_{\mathrm{vac}}},
\end{equation}
where $\Phi_{\mathrm{vac}}$ is vacuum flux or activity and $C_{\mathrm{vac}}$ is the effective transport capacity, stiffness, or constraint of the vacuum medium.

The question is whether apparent anomalies arise where vacuum flux and vacuum constraint are forced into extreme configurations. The answer is not asserted as established physics; it is formulated as a sequence of hypotheses with explicit failure criteria.

\begin{figure}[ht]
  \centering
  \includegraphics[width=0.95\textwidth]{figures/fig8_cosmology_bridge.pdf}
  \caption{Cosmological bridge map. Dark matter, dark energy, black-hole saturation, and the Blancken-layer proposal are shown as hypothesis classes attached to the same constrained-transport operator. The lower bar states the required falsification standard.}
  \label{fig:cosmology_bridge}
\end{figure}

\section{Vacuum Constraint Geometry}

The common mathematical object is the constrained transport operator:
\begin{equation}
    \mathcal{T}[\Phi,C] =
    \nabla\cdot\left(C^{-1}\nabla\Phi\right).
\end{equation}
In ordinary diffusion, transport depends only on gradient and conductivity. In CDFD, conductivity is itself a dynamical field. Large flux gradients alter constraint; altered constraint bends subsequent transport. This feedback is the mechanism by which geometry can emerge from transport.

\begin{hypothesis}[Vacuum constraint feedback]
The vacuum behaves, at an effective level, as a transport medium whose constraint field responds to persistent flux stress:
\begin{equation}
    \partial_t C_{\mathrm{vac}}
    =
    \alpha_{\mathrm{vac}}|\Phi_{\mathrm{vac}}|
    -\beta_{\mathrm{vac}}C_{\mathrm{vac}}
    +\gamma_{\mathrm{vac}}\nabla^2 C_{\mathrm{vac}}.
\end{equation}
\end{hypothesis}
Here $\alpha_{\mathrm{vac}}$ is the vacuum constraint-reinforcement coefficient,
$\beta_{\mathrm{vac}}$ is the vacuum relaxation coefficient, and
$\gamma_{\mathrm{vac}}$ is the spatial smoothing or diffusion coefficient for the
effective constraint field.

This is the bridge assumption. If no such effective feedback exists, the rest of this paper collapses into metaphor.

\section{Archive Source, Reproducibility, and Claim Discipline}

This manuscript expands the archived physical-mysteries synthesis rather than
discarding it. The archive contained five compact source claims:
\begin{enumerate}
    \item dark matter as geometric vacuum hardening;
    \item dark energy as adaptive metric expansion;
    \item black holes as transport-capacity rupture;
    \item the mass spectrum as vacuum topology;
    \item the Blancken layer as a pre-geometric substrate.
\end{enumerate}
Those claims are retained, but their status is corrected. The paper does not
claim to have replaced general relativity, particle dark matter searches, or
precision cosmology. It gives a constraint-transport research program whose
cosmological claims must produce independent discriminants before they can be
called physical explanations.

The original synthesis figures were archive assets. In this version they are
regenerated by \texttt{supplementary\_XI.py} into
\texttt{outputs/paper\_XI/}, with CSV tables and check files recording the
toy assumptions behind each conceptual panel. The figures remain schematics;
their value is that the bridge claims now have an auditable source trail rather
than orphaned images.

\begin{center}
\begin{tabular}{p{0.24\textwidth}p{0.28\textwidth}p{0.34\textwidth}}
\toprule
\textbf{Archive claim} & \textbf{Retained kernel} & \textbf{Correct status here}\\
\midrule
Dark matter & constraint gradients can mimic excess lensing & hypothesis requiring residual-lensing tests\\
Dark energy & vacuum self-consistency can leave pressure residuals & hypothesis linked to Paper X\\
Black holes & horizons are capacity saturation boundaries & interpretive bridge to GR\\
Mass topology & stable masses are topological constraint states & inherited from Papers I--VI\\
Blancken layer & pre-geometric substrate below spacetime & speculative until invariant is derived\\
\bottomrule
\end{tabular}
\end{center}

\section{Dark Matter as Constraint-Induced Lensing}

The CDFD dark-matter interpretation is not that ordinary matter is secretly missing from observations. It is that some portion of the gravitationally inferred anomaly could be modeled as vacuum constraint geometry induced by persistent flux aggregation.

Let $\rho_b$ be baryonic mass density and let $\Delta C_{\mathrm{vac}}$ be the induced constraint deformation. The effective lensing potential is written schematically:
\begin{equation}
    \Phi_{\mathrm{lens}}
    =
    \Phi_{\mathrm{GR}}[\rho_b]
    +
    \Phi_{\mathrm{CDFD}}[\Delta C_{\mathrm{vac}}].
\end{equation}
The second term is the hypothesis. It says that the vacuum remembers or stiffens around long-lived flux structures, producing lensing behavior not accounted for by visible matter alone.

\begin{prediction}[Geometry-memory signature]
If dark matter is partly vacuum-constraint geometry, lensing residuals should correlate not only with present baryonic mass but also with indicators of long-term flux history: merger pathways, rotational shear, plasma history, or persistent high-throughput structure.
\end{prediction}

\begin{criterion}[Dark-matter falsification]
The CDFD dark-matter claim fails if lensing residuals require no history-dependent or constraint-like term once baryons, standard dark matter models, and modified-gravity alternatives are properly fitted.
\end{criterion}

\begin{figure}[ht]
  \centering
  \includegraphics[width=0.82\textwidth]{figures/paper_4_dark_matter_lensing.png}
\caption{Reproducible Paper XI conceptual output for the dark-matter claim:
  persistent flux aggregation hardens the surrounding constraint geometry,
  producing lensing-like curvature. In this manuscript the image is treated
  as a conceptual model, not as observational evidence.}
  \label{fig:archive_dark_matter_lensing}
\end{figure}

\subsection{Observable Residual Protocol}

The dark-matter claim becomes scientific only if it specifies what to fit.
The proposed observational object is a residual field
\begin{equation}
    R_{\mathrm{lens}}(\mathbf{x})
    =
    \Phi_{\mathrm{obs}}(\mathbf{x})
    -
    \Phi_{\mathrm{GR+baryon}}(\mathbf{x})
    -
    \Phi_{\mathrm{standard\,DM}}(\mathbf{x}|\Theta),
\end{equation}
where $\Theta$ denotes the best-fit parameters of a conventional dark-matter
model. A CDFD contribution must then correlate with a measurable history
functional
\begin{equation}
    H_C(\mathbf{x},t)
    =
    \int_{-\infty}^{t}
    K(t-t')
    \left[
    \alpha_{\mathrm{vac}}|\Phi(\mathbf{x},t')|
    +\gamma_{\mathrm{vac}}\nabla^2C(\mathbf{x},t')
    \right]dt',
\end{equation}
not merely with current baryonic mass. The falsifiable statement is:
clusters or galaxies with similar baryonic mass but different merger,
plasma, or shear histories should show different residual constraint
patterns if the CDFD term is real.

\subsection{Relation to Modified Gravity}

CDFD is not identical to modified Newtonian dynamics or a generic
modified-gravity ansatz. In this paper the modification is path-dependent:
the constraint field has memory. A static algebraic relation between
acceleration and baryonic matter would not be sufficient to confirm the
CDFD interpretation. The distinguishing feature must be hysteresis,
relaxation, or spatial lag in the inferred constraint field.

\section{Dark Energy as Residual Vacuum Self-Consistency}

The vacuum equation-of-state paper~\cite{MujjabiX2026} develops the condition
$\Psi_{s,\mathrm{vac}}\to1$ as a stability point. In the cosmological bridge,
dark energy is interpreted as the residual pressure of a vacuum that approaches,
but has not globally completed, self-consistency.

At perfect local balance:
\begin{equation}
    \Phi_{\mathrm{vac}} \approx C_{\mathrm{vac}}.
\end{equation}
The residual term is:
\begin{equation}
    \delta_{\mathrm{vac}}
    =
    \frac{\Phi_{\mathrm{vac}}-C_{\mathrm{vac}}}{C_{\mathrm{vac}}}.
\end{equation}
If $\delta_{\mathrm{vac}}$ remains nonzero on cosmological scales, it may appear as a large-scale pressure term.

\begin{hypothesis}[Residual expansion]
Cosmic acceleration is a residual of incomplete vacuum self-consistency at Hubble scale rather than an independent substance.
\end{hypothesis}

\begin{prediction}[Scale dependence]
The effective dark-energy term should be close to constant where the vacuum is already equilibrated, but may show subtle scale or epoch dependence where $\nabla^2 C_{\mathrm{vac}}$ remains significant.
\end{prediction}

\begin{figure}[ht]
  \centering
  \includegraphics[width=0.82\textwidth]{figures/paper_4_dark_energy_expansion.png}
\caption{Reproducible Paper XI conceptual output for adaptive expansion. The
  revised claim is narrower: cosmic acceleration is modeled as a residual
  of incomplete vacuum self-consistency, not as proof that the residual has
  already been derived from first principles.}
  \label{fig:archive_dark_energy_expansion}
\end{figure}

\subsection{Connection to Paper X}

Paper X gives the local vacuum condition
\begin{equation}
    \Psi_{s,\mathrm{vac}}=\frac{\Phi_{\mathrm{vac}}}{C_{\mathrm{vac}}}\to1.
\end{equation}
The cosmological extension is not automatic. It requires a scale-separated
residual:
\begin{equation}
    \rho_{\Lambda,\mathrm{eff}}
    =
    \mathcal{R}\!\left[
      \Phi_{\mathrm{vac}}-C_{\mathrm{vac}},
      \nabla C_{\mathrm{vac}},
      \nabla^2 C_{\mathrm{vac}},
      H(t)
    \right],
\end{equation}
where $\mathcal{R}$ is the still-unfinished residual map. The archive
correctly identified the near-cancellation idea, but the mature manuscript
must state the open calculation: derive $\mathcal{R}$ with units, boundary
conditions, and cosmological parameters fixed before comparison with data.

\begin{figure}[ht]
  \centering
  \includegraphics[width=0.78\textwidth]{figures/fig10_1_panel.pdf}
  \caption{Archived Paper X vacuum equation-of-state check: the equilibrium variable $\Psi_s$ tracks $\Phi/C$ across the torus-knot family. In this synthesis paper the same relation is used only as a bridge assumption for cosmological hypotheses, not as direct evidence for the dark sector.}
  \label{fig:psi_phi_over_c_bridge}
\end{figure}

\begin{figure}[ht]
  \centering
  \includegraphics[width=0.78\textwidth]{figures/fig10_2_panel.pdf}
  \caption{Archived Paper X residual panel. Residuals near numerical precision support internal consistency of the vacuum equation-of-state calculation; external cosmological validation remains a separate requirement.}
  \label{fig:vacuum_residual_bridge}
\end{figure}

\section{Black Holes as Constraint Saturation Boundaries}

In CDFD language, a black hole is not merely high density. It is the limit where the transport system reaches a local boundary condition:
\begin{equation}
    J \rightarrow J_{\mathrm{crit}},
    \qquad
    \Psi_s \rightarrow \infty,
\end{equation}
because flux demand overwhelms outward transport capacity. The horizon is interpreted as the surface at which constraint geometry becomes causally dominant.

This does not remove the general relativistic horizon. It reframes the horizon as the macroscopic expression of a transport-capacity limit. The useful test is whether the CDFD variables add predictive structure to black-hole thermodynamics, ringdown deviations, accretion-state transitions, or information-flow constraints.

\begin{criterion}[Black-hole falsification]
The CDFD black-hole interpretation fails if it yields no measurable distinction from ordinary general relativity or from known effective-field corrections.
\end{criterion}

\begin{figure}[ht]
  \centering
  \includegraphics[width=0.82\textwidth]{figures/paper_4_black_hole_rupture.png}
\caption{Reproducible Paper XI conceptual output for black-hole rupture. The
  figure is retained as a transport-capacity visualization; the actual
  physics claim must be tested through horizon, ringdown, or thermodynamic
  discriminants.}
  \label{fig:archive_black_hole_rupture}
\end{figure}

\subsection{Capacity Boundary Formulation}

Let $J$ denote outward information/energy transport through a surface
$\mathcal{S}$ and let $J_{\mathrm{crit}}[C]$ denote the local capacity
allowed by the constraint geometry. A CDFD horizon is the boundary at which
\begin{equation}
    \frac{J}{J_{\mathrm{crit}}[C]}\rightarrow1^{-},
    \qquad
    \partial_r\left(\frac{J}{J_{\mathrm{crit}}[C]}\right)>0.
\end{equation}
This is not meant to erase the relativistic event horizon. It reframes the
horizon as a saturation surface in a transport medium. A useful next
calculation would show whether $J_{\mathrm{crit}}[C]$ reproduces the
Schwarzschild radius in the classical limit and whether it predicts any
controlled deviation in rotating or charged geometries.

\subsection{Where the Claim Can Fail}

The black-hole bridge fails if the capacity-boundary language cannot
recover ordinary GR in regimes where GR is already verified. It also fails
if it produces no additional observable beyond a relabeling of established
geometry. Candidate discriminants are ringdown mode shifts, accretion-state
hysteresis, entropy-area corrections, or information-flow bounds. Without
one of these, this section remains a conceptual translation.

\section{The Blancken Layer}

The Blancken layer is the most speculative part of the physics sequence. It is proposed as a pre-geometric substrate below ordinary spacetime description. In this view, spacetime is not fundamental but emerges as a stable constraint phase of a deeper continuous medium.

The role of the Blancken layer is conceptual: it provides a place for the continuous field before metric geometry, particle topology, or classical spacetime have appeared. A responsible formulation must keep it hypothesis-level until it produces mathematical consequences.

\begin{hypothesis}[Pre-geometric substrate]
There exists a sub-Planckian continuous layer whose stable macroscopic excitations appear as vacuum geometry, particle topology, and cosmic expansion.
\end{hypothesis}

\begin{criterion}[Blancken-layer falsification]
The Blancken-layer claim remains non-scientific unless it yields a calculable invariant, scaling law, or observational discriminant not already supplied by established theory.
\end{criterion}

\begin{figure}[ht]
  \centering
  \includegraphics[width=0.82\textwidth]{figures/paper_4_blancken_layer.png}
\caption{Reproducible Paper XI conceptual output for the Blancken layer. In the revised manuscript the
  layer is explicitly marked as speculative until it yields a calculable
  invariant or boundary condition.}
  \label{fig:archive_blancken_layer}
\end{figure}

\subsection{Minimum Mathematical Requirement}

The Blancken layer cannot remain a name for ``whatever is below spacetime.''
For the proposal to become physics, it must produce at least one of the
following:
\begin{enumerate}
    \item a dimensionless invariant preserved under emergence of metric
    geometry;
    \item a boundary condition that fixes the vacuum self-consistency
    selector used across Papers VI--X;
    \item a scaling law connecting sub-Planckian constraint to an observed
    constant or cosmological residual;
    \item a no-go theorem showing which effective geometries cannot emerge
    from constrained transport.
\end{enumerate}
The most valuable target is item 2: if the Blancken layer can derive the
fixed-point condition $n\theta_n=\widetilde{Q}_n(\theta_n)$, it would
strengthen both the torus-knot hierarchy and the cosmological bridge.

\section{Cosmological Mujjabi Tests}

The cosmology bridge is the large-scale version of the same question asked in
Paper~XII: does constraint history leave an observable residual after standard
models are fitted? Three targets follow from the source notes and the Part I
grammar:
\begin{enumerate}
    \item \textbf{Alpha drift as response drift.} Here un-subscripted $\alpha$
    denotes the electromagnetic fine-structure constant. If $S$ or $M_s$ varies over
    cosmological environments, high-precision spectral data should bound or
    reveal a correlated shift in effective $\alpha$.
    \item \textbf{Lensing hysteresis.} If dark-matter-like effects are
    constraint-history geometry, residual lensing should correlate with a
    history functional, not only with instantaneous baryonic mass.
    \item \textbf{Horizon information bounds.} If black holes are capacity
    saturation boundaries, ringdown, entropy-area, or information-flow
    observables must recover GR first and then state a controlled correction.
\end{enumerate}
These are Mujjabi Tests at cosmological scale. They are not evidence until they
produce a discriminant that standard cosmology cannot absorb with fewer
assumptions.

\section{Connection to the Mass Spectrum}

The physics series treats particle masses as self-consistent topological states of constrained vacuum transport. The dark-sector synthesis depends on the same premise: stable objects appear when continuous flux becomes locked by constraint topology. The difference is scale. In the lepton papers, the topology is microscopic and quantified through mass ratios and knot families. In the cosmological bridge, the topology is macroscopic and expressed through lensing, horizons, and expansion.

This scale analogy is not proof. It is a proposed research program:
\begin{enumerate}
    \item derive microscopic self-consistency conditions;
    \item test mass-spectrum predictions;
    \item extend the same constraint operator to vacuum geometry;
    \item search for dark-sector observables that require a constraint-history term.
\end{enumerate}

\section{Submission-Grade Falsification Program}

For this paper to stand beside the more mathematical physics papers, each
cosmological bridge must be tied to a concrete failure route:
\begin{center}
\begin{tabular}{p{0.22\textwidth}p{0.34\textwidth}p{0.30\textwidth}}
\toprule
\textbf{Bridge} & \textbf{Primary test} & \textbf{Failure condition}\\
\midrule
Dark matter & residual lensing vs history functional $H_C$ & no history-dependent residual after standard fits\\
Dark energy & derivation of $\rho_{\Lambda,\mathrm{eff}}$ from $\Psi_{s,\mathrm{vac}}\to1$ & wrong units, free tuning, or $w$ incompatible with data\\
Black holes & capacity-boundary recovery of GR plus controlled correction & no recovery of GR or no distinct prediction\\
Blancken layer & invariant or boundary condition below metric geometry & no calculable quantity beyond metaphor\\
Mass topology bridge & derivation of fixed-point selector from vacuum EOS & spectra need fitted selectors per family\\
\bottomrule
\end{tabular}
\end{center}

This is deliberately stricter than the archive language. The older draft
captured the conceptual map; this version states what would make the map
scientific and what would break it.

\section{Reproducibility Artifacts}

The accompanying script \texttt{supplementary\_XI.py} generates the figures
and tables used by this manuscript. The notebook
\texttt{notebooks/paper\_XI\_fullstack.ipynb} calls the same function for an
interactive reproduction path.

\begin{center}
\small
\begin{tabular}{p{0.42\textwidth}p{0.42\textwidth}}
\toprule
\textbf{Artifact} & \textbf{Role}\\
\midrule
\url{table1_cosmology_bridge_map.csv} & bridge taxonomy and failure criteria\\
\url{table2_dark_matter_residual_sample.csv} & toy residual sample used for the lensing-history gate\\
\url{table3_dark_energy_residual_curve.csv} & epoch-dependent residual-pressure proxy\\
\url{table4_black_hole_capacity_boundary.csv} & transport-saturation boundary proxy\\
\url{table5_blancken_layer_requirements.csv} & minimum mathematical requirements for the substrate claim\\
\url{checks_summary_XI_gate.csv} & consolidated reproducibility gate\\
\bottomrule
\end{tabular}
\end{center}

\section{Summary of Status}

\begin{center}
\begin{tabular}{p{0.25\textwidth}p{0.28\textwidth}p{0.32\textwidth}}
\toprule
\textbf{Claim} & \textbf{Current status} & \textbf{Required next step}\\
\midrule
Vacuum $\Psi_s=1$ stability & derived within CDFD assumptions & independent mathematical review\\
Lepton mass topology & model-derived and numerically checkable & precision comparison and external critique\\
Dark matter as constraint geometry & hypothesis & lensing residual discriminant\\
Dark energy as residual self-consistency & hypothesis & scale/epoch-dependent prediction\\
Black holes as transport saturation & hypothesis & ringdown or thermodynamic discriminant\\
Blancken layer & speculative hypothesis & calculable invariant\\
\bottomrule
\end{tabular}
\end{center}

\section{Conclusion}

This synthesis is the edge of the physics program, not its foundation. The strongest CDFD physics claims remain the self-consistency and mass-spectrum derivations developed elsewhere in the series. The dark matter, dark energy, black-hole, and Blancken-layer claims are valuable only if they mature into concrete discriminants. Until then, they should be presented as a structured hypothesis: one transport ontology, several cosmological anomalies, and explicit routes by which the whole synthesis can fail.

\section*{Supplementary Material}

This manuscript is accompanied by release-archive supplementary material in
the canonical Part I tree. The relevant paper-local Python scripts, numerical checks, generated figure outputs,
tabulated data, figure assets, and environment notes are maintained under
\texttt{notebooks/}, \texttt{outputs/}, \texttt{figures/}, and
\texttt{REPRODUCIBILITY.md} in \texttt{Part\_I\_Fundamental\_Physics/}.


\section*{Data and Code Availability}

The release-archive supplementary material is included in the archived Part I
release on Zenodo at \url{https://doi.org/10.5281/zenodo.20250820}. The current version DOI is assigned by Zenodo after each GitHub release. Zenodo is the
permanent release archive, and the public source archive is available on GitHub at
\url{https://github.com/bampita-bico/CDFD-Part-I-Release}. The broader public CDFD
Runtime is separate reusable software infrastructure; the numerical values reported
here are reproduced from this paper's Supplementary Material.

\section*{Acknowledgments}

The author used AI-assisted tools for formatting checks and code-verification
support. The scientific claims, interpretations, and final manuscript decisions
are the author's responsibility.

\section*{Declarations}

\subsection*{Funding}
This research received no external institutional funding.

\subsection*{Author Contributions}
Steve Bico Mujjabi: conceptualization, methodology, formal analysis, software,
validation, visualization, writing--original draft, and writing--review and
editing.

\subsection*{Conflicts of Interest}
The author declares no conflicts of interest.

\bibliographystyle{apsrev4-2}
\bibliography{references}

\end{document}
